\clearpage
\fchapter{Objeto}


El objeto del presente proyecto es la integración, configuración y puesta en marcha de un sistema de bin picking inteligente basado en la plataforma SIMATIC Robot Pick AI de Siemens, combinando robótica colaborativa, visión artificial 3D e inteligencia artificial aplicada a la detección y manipulación de piezas en entornos no estructurados. El sistema resultante constituye una solución completa de automatización flexible, capaz de operar sin necesidad de posicionamiento previo de las piezas.
Para ello me propongo alcanzar los siguientes objetivos:

\begin{itemize}
	\item Configurar y verificar la comunicación entre todos los elementos del sistema (cámara Orbbec Femto Mega, IPC, PLC S7-1500 y robot UR3e) mediante protocolo Profinet en la red 192.168.15.xxx.

	\item Realizar la configuración del software PickAI en el IPC, incluyendo la calibración de la cámara, la calibración del robot y la configuración de la herramienta de agarre.

	\item Adaptar el programa de TIA Portal proporcionado en la nota de aplicación de Siemens, ajustando los parámetros del DB4 (dirección IP del IPC y medidas del contenedor) a las características del equipamiento disponible.

	\item Adaptar el programa del robot UR3e, ajustando los puntos de paso y las posiciones de depósito al entorno de trabajo real.

	\item Validar el funcionamiento completo del sistema realizando pruebas de detección, agarre y depósito de piezas, verificando que el robot opera de forma autónoma hasta vaciar el contenedor.
\end{itemize}

La identificación individual de las piezas para clasificarlas (mediante RFID o códigos QR) queda fuera del alcance de este proyecto y se plantea como ampliación futura en el apartado de posibilidades de innovación y mejoras.
